\section{Pullback and descent}
\label{sec:psh-descent}

Pullback preserves the local positivity of a potential. Compact fibers force descent when the background has no vertical component; under a modification, a Jacobian factor accounts for the change in multiplier-ideal integrability.

\Needspace{18\baselineskip}
\subsection{Holomorphic disks}

A smooth background need not have local potentials, but its pullback to a disk does. Testing there retains the exact lower curvature bound under a holomorphic map.

\begin{lemma}\label{lem:smooth-disk-criterion}
Let $\theta$ be a smooth real $(1,1)$-form on a complex manifold $M$, and let $u$ be quasi-plurisubharmonic on $M$. Then $\theta+\ddc u\geq0$ if and only if the following condition holds. For every holomorphic map $f:\Delta\to M$ from the unit disk whose pullback $u\circ f$ is not identically $-\infty$, the function $u\circ f+a$ is subharmonic wherever $a$ is a smooth local potential for $f^*\theta$.
\end{lemma}
\begin{proof}
Suppose first that $\theta+\ddc u\geq0$. On a coordinate ball convolve with a smooth nonnegative radial kernel. The smooth functions $u_\epsilon$ satisfy
\[
 \ddc u_\epsilon\geq-\theta_\epsilon,
 \qquad\theta_\epsilon\longrightarrow\theta\quad\text{smoothly}.
\]
Write $u=p-r$, with $p$ plurisubharmonic and $r$ smooth. The convolutions $p_\epsilon$ decrease to $p$. On a smaller disk where $u\circ f$ is genuine, $p_\epsilon\circ f$ therefore converges in $L^1_{\mathrm{loc}}$ to $p\circ f$, while $r_\epsilon\circ f$ converges smoothly. Pull back the smooth inequality and pass to distributions. This gives $\ddc(u\circ f)+f^*\theta\geq0$. Since $u\circ f=p\circ f-r\circ f$ is quasi-subharmonic, the distributional inequality makes $u\circ f+a$ subharmonic.

Conversely, test on affine complex lines in a coordinate ball. For any fixed nonzero constant direction $\xi$, local integrability of $u$ implies that almost every parallel line has a genuine restriction. Fubini's theorem applied to nonnegative smooth test functions gives positivity of the scalar distribution obtained by evaluating the Hermitian coefficient matrix of $\theta+\ddc u$ on $\xi$. Positivity for a countable dense set of constant directions proves positivity of the Hermitian matrix of coefficient measures, hence the current inequality.
\end{proof}

Consequently, holomorphic pullback preserves the background inequality wherever the pulled-back potential is genuine: apply the disk criterion to compositions of maps. We also use descent of subharmonicity under a finite disk map $\rho(t)=t^k$, $k\geq1$. If $h\circ\rho$ is subharmonic, local inverse branches make $h$ subharmonic off zero. It is bounded above near zero and has the canonical extension value there, since
\[
 h(0)=(h\circ\rho)(0)
      =\limsup_{t\to0,\ t\ne0}h(t^k)
      =\limsup_{z\to0,\ z\ne0}h(z).
\]
The removable-singularity theorem makes $h$ subharmonic at zero as well.

\Needspace{17\baselineskip}
\subsection{Compact fibers}

The connected components of a submersion fiber determine a finite covering of the base. Potentials first descend to that covering.

\begin{lemma}\label{lem:submersion-component-cover}\index{Descent!cover of fiber components}\index{Submersion!cover of fiber components}
Let $p:M\to Y$ be a surjective holomorphic submersion between connected compact complex manifolds. There is a factorization
\[
 M\xrightarrow{\pi}Y'\xrightarrow{r}Y,\qquad p=r\circ\pi,
\]
where $Y'$ is a connected compact complex manifold, $r$ is a finite unramified holomorphic covering and $\pi$ is a surjective holomorphic submersion with connected fibers. The points of $r^{-1}(y)$ correspond to the connected components of $p^{-1}(y)$, which are the fibers of $\pi$ over those points.
\end{lemma}
\begin{proof}
Fix a smooth Hermitian metric on $M$. The orthogonal complement of $\ker\mathrm dp$ gives a smooth horizontal lift of tangent vectors on $Y$. Around $y_0\in Y$ choose a coordinate ball $B$ centered at $y_0$, with compact closure inside a larger coordinate ball. For $y\in B$, lift the straight segment from $y_0$ to $y$. The ordinary differential equation for horizontal transport has a solution along the whole segment: its base remains in the compact closure of $B$, whose inverse image is compact. Smooth dependence on the starting point and on $y$, together with transport along the reversed segment, gives a smooth trivialization
\[
 p^{-1}(B)\simeq B\times p^{-1}(y_0)
\]
over $B$.

The compact smooth manifold $p^{-1}(y_0)$ has finitely many connected components. The trivialization partitions $p^{-1}(B)$ into finitely many open and closed pieces, each with one connected component in every fiber. On overlaps, two such local partitions identify the same fiber components; their labels differ by locally constant permutations. Glue one copy of each base ball for every component using these identifications. Equivalently, take the quotient of $M$ that identifies points in the same connected component of a fiber. The local partitions give this quotient a Hausdorff covering topology: distinct points over one base point are separated by different sheets, and points over different base points are separated downstairs.

Call the quotient $Y'$ and its projection $r$. Base coordinates on each sheet give holomorphic charts, with holomorphic transition maps, so $r$ is a local biholomorphism. It is a finite covering and $Y'$ is compact, since a finite cover of a compact base is compact. The quotient map $\pi$ is locally $p$ restricted to one of the open pieces. It is therefore a holomorphic submersion, with the asserted connected fibers. Surjectivity of $\pi$ and connectedness of $M$ make $Y'$ connected.
\end{proof}

\Needspace{17\baselineskip}
\begin{proposition}\label{prop:psh-descent-submersion}\index{Descent!submersion}
Let $\pi:Y\to X$ be a surjective holomorphic submersion between connected compact complex manifolds, with connected fibers. Let $\theta$ be a smooth real $(1,1)$-form on $X$. Pullback gives a bijection
\[
 \PSH(X,\theta)\longrightarrow\PSH(Y,\pi^*\theta).
\]
Every member on the right is constant on each fiber, with the value $-\infty$ allowed on an entire fiber.
\end{proposition}
\begin{proof}
In a coordinate ball $V\subset X$, choose a smooth function $\rho$ with $\ddc\rho\geq\theta$. If $u\in\PSH(Y,\pi^*\theta)$, then $u+\rho\circ\pi$ is plurisubharmonic on $\pi^{-1}(V)$. Its restriction to a fiber is plurisubharmonic or identically $-\infty$. Since a fiber is compact and connected, the maximum principle makes it constant in either case. Thus $u=v\circ\pi$ for a uniquely defined function $v$ on $X$.

A holomorphic submersion admits local holomorphic sections. Composing $u+\rho\circ\pi$ with such a section shows that $v+\rho$ is plurisubharmonic or identically $-\infty$ on a smaller ball. The latter alternative on a nonempty open ball would make $u=-\infty$ on its open inverse image, contradicting genuineness. Hence $v$ is quasi-plurisubharmonic. Distributional differentiation in local product coordinates gives
\[
 \pi^*(\theta+\ddc v)=\pi^*\theta+\ddc u\geq0.
\]
Positivity of a current of this pulled-back form implies positivity on the base: test in product coordinates and integrate a nonnegative compactly supported volume form of integral one in the fiber variables. It follows that $v\in\PSH(X,\theta)$.

Conversely, \cref{lem:smooth-disk-criterion} shows that pullback of a $\theta$-plurisubharmonic function by a submersion is $\pi^*\theta$-plurisubharmonic. It is genuine by surjectivity and local product coordinates. Uniqueness follows from surjectivity.
\end{proof}

\Needspace{11\baselineskip}
\subsection{Modifications}

\begin{proposition}\label{prop:psh-descent-modification}\index{Descent!modification}
Let $\pi:Y\to X$ be a modification of connected compact complex manifolds, and let $\theta$ be a smooth real $(1,1)$-form on $X$. Pullback gives a bijection
\[
 \PSH(X,\theta)\longrightarrow\PSH(Y,\pi^*\theta).
\]
\end{proposition}
\begin{proof}
Let $S\subset X$ be a proper analytic set over whose complement $\pi$ is an isomorphism, and put $D=\pi^{-1}(S)$. Pullback preserves the curvature inequality by \cref{lem:smooth-disk-criterion} and is genuine on the isomorphism locus. It is injective because $\pi$ is surjective.

To prove surjectivity, let $b\in\PSH(Y,\pi^*\theta)$ and define $u$ off $S$ by the inverse map. On a coordinate neighborhood choose a smooth $\rho$ with $\ddc\rho\geq\theta$. Then $b+\rho\circ\pi$ is plurisubharmonic. Properness and its local upper bounds give a local upper bound for $u+\rho$ near $S$, so the removable-singularity theorem extends $u+\rho$ plurisubharmonically across $S$. The resulting quasi-plurisubharmonic extensions of $u$ agree almost everywhere on overlaps and hence glue as canonical representatives.

The current $\theta+\ddc u$ is positive off $S$. Locally it is the sum of the positive current $\ddc(u+\rho)$ and the smooth form $\theta-\ddc\rho$. The latter has zero coefficient-measure mass on $S$, so the restriction of $\theta+\ddc u$ to $S$ is positive as well. Thus $u\in\PSH(X,\theta)$. The plurisubharmonic functions $b+\rho\circ\pi$ and $(u+\rho)\circ\pi$ agree off $D$, hence almost everywhere and everywhere. Therefore $b=u\circ\pi$.
\end{proof}

A modification preserves $L^1$ convergence within a family with one lower curvature bound. Compactness controls the pullbacks near the exceptional locus, where a direct change-of-variables estimate would involve a vanishing Jacobian.

\begin{lemma}\label{lem:modification-lone-pullback}
Let $\pi:Y\to X$ be a modification between connected complex manifolds, and let $\theta$ be a smooth real $(1,1)$-form on $X$. If $v_j\to v$ in $L^1_{\mathrm{loc}}(X)$ and all $v_j,v$ belong to $\PSH(X,\theta)$, then
\[
 v_j\circ\pi\longrightarrow v\circ\pi
       \quad\hbox{in }L^1_{\mathrm{loc}}(Y).
\]
In particular, on compact manifolds this is convergence in $L^1$.
\end{lemma}
\begin{proof}
\Cref{lem:qpsh-hartogs} bounds $v_j$ uniformly above on relatively compact subsets of $X$. The pullbacks are therefore locally uniformly bounded above and have the common smooth lower curvature bound $-\pi^*\theta$. After a smooth local correction, plurisubharmonic compactness gives a convergent subsequence on each connected coordinate region in $Y$, or convergence to $-\infty$ locally uniformly.

Every such region meets the complement of the exceptional locus. There $\pi$ is a biholomorphism, and change of variables gives $L^1_{\mathrm{loc}}$ convergence to $v\circ\pi$. Thus collapse is excluded. Each subsequential limit agrees with $v\circ\pi$ off the exceptional locus, hence almost everywhere, since that analytic locus has smooth volume zero. The limit is unique. Applying this argument to every subsequence and a countable coordinate cover proves convergence of the whole sequence in $L^1_{\mathrm{loc}}(Y)$.
\end{proof}

Capacity convergence is also preserved by a modification. The comparison takes place away from small neighborhoods of the exceptional sets; a uniform comparison of capacities on the whole manifolds is unnecessary.

\Needspace{10\baselineskip}
\begin{lemma}\label{lem:modification-capacity-pullback}
Let $\pi:Y\to X$ be a modification between connected compact complex manifolds, equipped with Hermitian forms $\omega_Y,\omega_X$. Write $\operatorname{Cap}_Y=\operatorname{Cap}_{\omega_Y}$ and $\operatorname{Cap}_X=\operatorname{Cap}_{\omega_X}$. For quasi-plurisubharmonic functions $v_j,v$ on $X$,
\[
 v_j\longrightarrow v\text{ in capacity on }X
 \quad\Longleftrightarrow\quad
 v_j\circ\pi\longrightarrow v\circ\pi\text{ in capacity on }Y.
\]
\end{lemma}
\begin{proof}
Choose a proper analytic set $S\subset X$ such that $\pi$ is an isomorphism outside $D=\pi^{-1}(S)$. Both sets have capacity zero. For $\epsilon>0$, choose an open neighborhood $V$ of $D$ of capacity less than $\epsilon$. Properness gives a neighborhood $U_0$ of $S$ with $\pi^{-1}(U_0)\subset V$. Shrink it to an open $U$ with $\operatorname{Cap}_X(U)<\epsilon$. The inverse map is biholomorphic near the compact set $X\setminus U$. A finite cover by nested coordinate charts and the local-global capacity comparison of \cref{sec:hermitian-capacity} give a constant $C_\epsilon$ such that, for every Borel set $E\subset X$,
\[
 \begin{split}
 \operatorname{Cap}_Y(\pi^{-1}E)&\leq\epsilon+C_\epsilon\operatorname{Cap}_X(E),\\
 \operatorname{Cap}_X(E)&\leq\epsilon+C_\epsilon\operatorname{Cap}_Y(\pi^{-1}E).
 \end{split}
\]
Apply these inequalities to $E=\{|v_j-v|>\delta\}$ for each fixed $\delta>0$, omitting the pluripolar pole sets. First let $j\to\infty$, then $\epsilon\downarrow0$.
\end{proof}

\Needspace{13\baselineskip}
\subsection{Multiplier ideals under modifications}
\label{sec:psh-birational-ideals}

Let $\pi:Y\to X$ be a modification of connected compact complex manifolds of dimension $n$. There is an analytic subset $S\subset X$ of codimension at least two such that $\pi$ is an isomorphism over $X\setminus S$. Let $K_X$ and $K_Y$ be the canonical holomorphic line bundles. Set $D=\pi^{-1}(S)$ and
\[
 R_\pi=K_Y\otimes\pi^*K_X^{-1}.
\]
The determinant of the differential is a canonical holomorphic section $J_\pi$ of $R_\pi$. In local coordinates its coefficient is the complex Jacobian determinant. Its zeros are contained in $D$. These standard facts about modifications of smooth manifolds include the case $n=1$, when the modification is an isomorphism.

\Needspace{20\baselineskip}
\begin{lemma}\label{lem:adjoint-ideal-transform}\index{Multiplier ideal!birational transform}
Let $a$ be quasi-plurisubharmonic on $X$. Under the natural identification away from $S$,
\[
 \pi_*\bigl(R_\pi\otimes\mathcal I(a\circ\pi)\bigr)
       =\mathcal I(a).
\]
More precisely, for any holomorphic line bundle $M$ on $X$, the map
\[
 s\longmapsto J_\pi\,\pi^*s
\]
identifies
\[
 H^0\bigl(X,M\otimes\mathcal I(a)\bigr)
 \simeq H^0\bigl(Y,\pi^*M\otimes R_\pi
                         \otimes\mathcal I(a\circ\pi)\bigr).
\]
\end{lemma}
\begin{proof}
The assertion is local on $X$. In coordinates and local frames the proposed map sends a holomorphic coefficient $f$ to $(f\circ\pi)\operatorname{Jac}\pi$. Change of variables gives the local integrability equivalence between
\[
 |f|^2\mathrm{e}^{-a}
 \quad\hbox{and}\quad
 |f\circ\pi|^2|\operatorname{Jac}\pi|^2\mathrm{e}^{-a\circ\pi}.
\]
For the implication from integrability upstairs to integrability near a downstairs point, take a relatively compact smaller neighborhood. Its inverse image has compact closure by properness, so finitely many integrability charts suffice. Smooth positive coordinate densities are comparable on these smaller charts. Analytic exceptional sets have smooth volume zero and do not change the integrals.

It remains to justify that every section upstairs has the required form. Divide it by $J_\pi$ off $D$ and use the inverse of $\pi$ there. This gives a holomorphic section of $M$ off $S$. Hartogs extension across the codimension-at-least-two analytic set $S$ supplies a holomorphic section $s$ on $X$. Its image $J_\pi\pi^*s$ agrees with the given section off $D$, hence everywhere by the identity theorem. The integrability equivalence then places $s$ in the required multiplier ideal. The same argument on each open set proves the sheaf identity and the global identification.
\end{proof}

In particular, the formula does not identify $\pi^*\mathcal I(a)$ with $\mathcal I(a\circ\pi)$. The Jacobian is part of the holomorphic coefficient upstairs, rather than a factor that may be deleted from the integrability test.

We also need invariance of the ideal \emph{order}, which asks for all positive weights. We give the local argument to avoid assuming that a modification of a small open subset of $Y$ extends over $X$.

\Needspace{10\baselineskip}
\begin{lemma}\label{lem:birational-ideal-order}
For quasi-plurisubharmonic $a,b$ on $X$,
\[
 a\preceq_{\mathcal I}b
 \quad\Longleftrightarrow\quad
 a\circ\pi\preceq_{\mathcal I}b\circ\pi.
\]
The assertion is local on $X$ and also holds for a proper bimeromorphic map between noncompact smooth complex manifolds.
\end{lemma}
\begin{proof}
The reverse implication follows at every positive weight from the sheaf identity in \cref{lem:adjoint-ideal-transform}: tensor the upstairs ideal inclusion with $R_\pi$ and apply the left-exact direct image functor. The proof of that sheaf identity is local on the target; properness makes each fiber compact and supplies the required finite integrability cover. Compactness of the entire manifolds is unnecessary.

For the forward implication, fix a coordinate neighborhood $W\subset Y$, a smooth local modification $\sigma:Z\to W$, and a prime divisor $E$ on $Z$. Shrinking $W$ as necessary near the center under consideration, its image is contained in a coordinate neighborhood $U\subset X$. Put $F=\pi\circ\sigma$ and
\[
 A_F(E)=1+\operatorname{ord}_E\operatorname{Jac}F.
\]
The map $F$ is holomorphic and is one-to-one and biholomorphic off a proper analytic set. It need not be a proper modification of $U$.

For a positive integer $m$, let $\nu_E$ denote the generic Lelong number along $E$, and let $\operatorname{ord}_E$ denote the generic order of an ideal there. The proof of \cref{lem:bergman-divisorial-error} applies to this map on relatively compact coordinate patches and gives
\[
 \max\left\{\nu_E(a\circ F)-\frac{A_F(E)}m,0\right\}
 \leq \frac1m\operatorname{ord}_E\bigl(\mathcal I(ma)\mathcal O_Z\bigr)
 \leq\nu_E(a\circ F).
\]
For clarity, the upper inequality pulls back the local Bergman lower bound on $U$. For the lower one, a germ $f\in\mathcal I(ma)$ pulls back to an integrable coefficient $(f\circ F)\operatorname{Jac}F$: on a relatively compact upstairs patch the change-of-variables integral is bounded by the integral over its downstairs image. Generic divisorial removal and the one-variable integral give
\[
 \operatorname{ord}_E(f\circ F)+A_F(E)>m\nu_E(a\circ F).
\]
Only local injectivity off the exceptional set and the compact patch are needed for this estimate, not properness of $F$ over $U$. The same estimates hold for $b$.

The hypothesis implies $\mathcal I(ma)\subseteq\mathcal I(mb)$ for every $m>0$. Orders of the pulled-back ideals therefore satisfy the reverse inequality. Letting $m\to\infty$ in the two estimates proves
\[
 \nu_E(a\circ F)\geq\nu_E(b\circ F).
\]
This holds for every local divisor over $Y$. Blowing up a point on any local modification recovers its pointwise Lelong number; in dimension one the point already is a divisor. The exact local criterion in \cref{thm:local-i-inputs} now proves the upstairs ideal order. No extension of $\sigma$ outside its local domain was used.
\end{proof}

\Needspace{13\baselineskip}
\begin{lemma}\label{lem:ideal-equivalence-modification}
For a proper bimeromorphic map $\rho:Z\to Y$ between smooth complex manifolds and genuine quasi-plurisubharmonic functions $a,b$ on $Y$,
\[
 a\sim_{\mathcal I}b
       \quad\Longleftrightarrow\quad
 a\circ\rho\sim_{\mathcal I}b\circ\rho.
\]
The statement is local on $Y$.
\end{lemma}
\begin{proof}
Apply the local order equivalence of \cref{lem:birational-ideal-order} to both $(a,b)$ and $(b,a)$. The two order inequalities are exactly equality of the multiplier ideals at every positive weight.
\end{proof}

\Needspace{20\baselineskip}
\subsection{Restriction of analytic approximants}

A quasi-plurisubharmonic function with zero generic Lelong number along a submanifold may still restrict to the function identically equal to $-\infty$ (\cref{sec:analytic-types}). Analytic regularization removes this obstruction: the restrictions of the analytic approximants are quasi-plurisubharmonic, and a fixed lower curvature bound controls their generic Lelong numbers along divisors on every fixed modification.

\begin{corollary}\label{cor:restricted-analytic-divisorial-bound}
Let $\theta$ be a smooth real $(1,1)$-form and $u\in\PSH(X,\theta)$, and let $u_m$, with curvature errors $\epsilon_m$, be the approximants of \cref{cor:nonclosed-analytic-approximation}. Let $Y\subset X$ be a connected compact complex submanifold along which the generic Lelong number of $u$ vanishes:
\[
 \nu(u,Y)=\inf_{y\in Y}\nu(u,y)=0.
\]
Then each restriction $u_m|_Y$ is quasi-plurisubharmonic on $Y$; in particular, it is not identically $-\infty$. If $\pi\colon Z\to Y$ is a modification with smooth source and $E\subset Z$ is a prime divisor, then
\[
 \sup_m\nu\bigl((u_m|_Y)\circ\pi,E\bigr)<\infty.
\]
\end{corollary}
\begin{proof}
For each $m$, choose $y\in Y$ with $\nu(u,y)<1/m$. Skoda's local integrability criterion gives $\mathcal I(mu)_y=\mathcal O_{X,y}$, so the image ideal $\mathcal I(mu)\mathcal O_Y$ does not vanish identically. Since $u_m$ has analytic type $(1/m,\mathcal I(mu))$, \cref{prop:analytic-type-calculus} shows that $u_m|_Y$ is quasi-plurisubharmonic on $Y$. Its pullback to $Z$ is quasi-plurisubharmonic as well, since $\pi$ is biholomorphic off a proper analytic subset.

\Needspace{6\baselineskip}
Restricting and pulling back the curvature inequality of \cref{cor:nonclosed-analytic-approximation}, which is allowed by \cref{lem:smooth-disk-criterion}, gives
\[
 \pi^*(\theta|_Y)+\ddc\bigl((u_m|_Y)\circ\pi\bigr)
 \geq-\epsilon_m\pi^*(\omega|_Y).
\]
Since the errors $\epsilon_m$ are bounded, a sufficiently large multiple of a single Hermitian form on $Z$ dominates $\pi^*(\theta|_Y)+\epsilon_m\pi^*(\omega|_Y)$ for all $m$. \Cref{lem:gauduchon-uniform-lelong}, applied on $Z$ in this fixed background, therefore bounds the generic Lelong numbers along $E$ uniformly in $m$.
\end{proof}

